%==============================================================================
% Bernard Benson, PhD --- Long-form Academic CV
% Compile with: pdflatex CV_BernardBenson.tex  (run twice for hyperref refs)
%           or: latexmk -pdf CV_BernardBenson.tex
%==============================================================================
\documentclass[11pt,letterpaper]{article}

\usepackage[T1]{fontenc}
\usepackage[utf8]{inputenc}
\usepackage{lmodern}
\usepackage[margin=1in]{geometry}
\usepackage{titlesec}
\usepackage{enumitem}
\usepackage{xcolor}
\usepackage{tabularx}
\usepackage[hidelinks]{hyperref}

\hypersetup{colorlinks=true, urlcolor=black, linkcolor=black}

% --- Section heading: small-caps with a full-width rule underneath ----------
\titleformat{\section}
  {\large\scshape\bfseries}
  {}{0em}{}[\vspace{-0.5ex}\titlerule]
\titlespacing{\section}{0pt}{1.4ex plus 0.4ex minus 0.2ex}{1.0ex}

\setlist[itemize]{leftmargin=1.4em, topsep=0.2ex, itemsep=0.2ex, parsep=0pt}

\setlength{\parindent}{0pt}
\pagestyle{empty}

% --- Helper: an experience/appointment entry --------------------------------
% \cventry{Role}{Organization, Location}{Dates}
\newcommand{\cventry}[3]{%
  \par\vspace{0.6ex}%
  \noindent\begin{tabularx}{\textwidth}{@{}X r@{}}
    \textbf{#1} & \textbf{#3}\\
    \textit{#2} & \\
  \end{tabularx}%
  \vspace{-0.4ex}%
}

% --- Helper: a simple two-column line (left text, right-aligned date) --------
\newcommand{\daterow}[2]{%
  \noindent\begin{tabularx}{\textwidth}{@{}X r@{}}#1 & #2\end{tabularx}\par
}

% --- Helper: a publication item; bolds the author name automatically ---------
\newcommand{\me}{\textbf{B.~Benson}}
\newlist{pubs}{enumerate}{1}
\setlist[pubs]{leftmargin=2.0em, label=[\arabic*], topsep=0.4ex, itemsep=0.6ex, parsep=0pt, labelsep=0.5em}

%==============================================================================
\begin{document}

% ----------------------------- HEADER ---------------------------------------
\begin{center}
  {\LARGE\bfseries Bernard Benson, PhD}\\[0.6ex]
  {\small
    \href{mailto:bernardvbenson@gmail.com}{bernardvbenson@gmail.com}
    \;$|$\; 803-931-2777
    \;$|$\; \href{https://www.linkedin.com/in/bernardbenson}{LinkedIn}
    \;$|$\; \href{https://scholar.google.com/citations?user=uiHWfLEAAAAJ}{Google Scholar}
    \;$|$\; \href{https://github.com/bernardbenson}{GitHub}
  }
\end{center}
\vspace{0.4ex}

% ------------------------- RESEARCH INTERESTS -------------------------------
\section{Research Interests}
Machine learning and deep learning for the physical sciences; heliophysics and space-weather
forecasting; time-series modeling and uncertainty quantification; natural language processing
and large language models (retrieval-augmented generation, agentic and MCP-based tools);
computer vision.

% ------------------------------ EDUCATION -----------------------------------
\section{Education}
\cventry{Ph.D., Electrical and Computer Engineering}{University of Alabama in Huntsville, Huntsville, AL}{2016 -- 2021}

\cventry{M.S., Electrical and Computer Engineering}{University of Alabama in Huntsville, Huntsville, AL}{2011 -- 2012}

\cventry{B.Tech., Electrical and Electronics Engineering}{Geethanjali College of Engineering and Technology}{2006 -- 2010}

% ----------------------------- APPOINTMENTS ---------------------------------
\section{Appointments \& Professional Experience}

\cventry{Research Scientist IV}{NASA ODSI, University of Alabama in Huntsville, Huntsville, AL}{May 2025 -- Present}
\begin{itemize}
  \item Developed ensemble large language models for multi-label text classification of scientific documents.
  \item Built a search wrapper modernizing the backend of NASA's Science Discovery Engine (SDE) using OpenSearch.
  \item Researched and developed MCP tools for PDS4 and SDE data-search agents.
\end{itemize}

\cventry{Generative AI Solutions Architect}{Megan Soft VA, Remote}{Dec 2024 -- May 2025}
\begin{itemize}
  \item Developed an evaluation framework for generative-AI chatbots and RAG systems using DeepEval and Ragas.
  \item Established KPIs and testing procedures to optimize performance and user satisfaction.
  \item Designed multi-chatbot orchestration strategies for agents with overlapping capabilities.
  \item Designed and implemented recommender systems for pharmacy benefits management.
\end{itemize}

\cventry{Data Scientist}{McLeod Software, Birmingham, AL}{Aug 2021 -- Nov 2024}
\begin{itemize}
  \item Led development and deployment of ML models predicting freight rates for over 1,100 trucking organizations, improving pricing accuracy and efficiency.
  \item Designed and deployed retrieval-augmented-generation (RAG) chatbots for customer support, reducing resolution times.
  \item Built and deployed LayoutLMv2-based models to automate document classification.
  \item Led research and analysis of AI applications across the transportation and logistics industry.
  \item Represented McLeod Software as a conference speaker at industry events on AI in transportation and logistics.
\end{itemize}

\cventry{Applied ML Researcher}{Frontier Development Lab (NASA / SETI Institute), Remote}{Jun 2021 -- Aug 2021}
\begin{itemize}
  \item Applied machine learning to model the effect of solar drag on satellite orbits as part of a NASA-sponsored research project.
  \item Developed a multivariate time-series model based on deep neural network ensembles to forecast space-weather indices.
  \item Earned a team excellence award for the project.
\end{itemize}

\cventry{Graduate Teaching Assistant}{University of Alabama in Huntsville, Huntsville, AL}{Aug 2016 -- May 2021}
\begin{itemize}
  \item Taught and supervised lab sections for CPE211L Introduction to Programming (C++).
  \item Taught EE316 Electrical Circuits and Electrical Design Lab (Summer 2017).
\end{itemize}

\cventry{RF Engineer}{Nexgen Wireless / Global Telecom Associates, Milwaukee, WI}{Mar 2013 -- Mar 2016}
\begin{itemize}
  \item Analyzed drive-test data and network logs to identify and resolve coverage gaps and interference issues, improving network performance.
  \item Visualized network performance with heatmaps and coverage plots to support data-driven optimization.
\end{itemize}

% ----------------------------- PUBLICATIONS ---------------------------------
\section{Publications}

\textbf{\itshape Refereed Journal Articles}
\begin{pubs}
  \item \me, W.~D.~Pan, A.~Prasad, G.~A.~Gary, and Q.~Hu, ``Forecasting Solar Cycle 25 Using Deep Neural Networks,'' \textit{Solar Physics}, vol.~295, no.~5, art.~65, 2020.
  \item T.~Singh, \me, S.~A.~Z.~Raza, T.~K.~Kim, N.~V.~Pogorelov, W.~P.~Smith, and C.~N.~Arge, ``Improving the Arrival-Time Estimates of Coronal Mass Ejections by Using Magnetohydrodynamic Ensemble Modeling, Heliospheric Imager Data, and Machine Learning,'' \textit{The Astrophysical Journal}, vol.~948, no.~2, art.~78, 2023.
  \item \me, W.~D.~Pan, G.~A.~Gary, Q.~Hu, and T.~Staudinger, ``Determining the Parameter for the Linear Force-Free Magnetic Field Model with Multi-Dipolar Configurations Using Deep Neural Networks,'' \textit{Astronomy and Computing}, vol.~26, pp.~50--60, 2019.
  \item T.~S.~Newman, C.~W.~Hall, L.~Farris, T.~Singh, N.~Pogorelov, \me, S.~Raza, \textit{et al.}, ``Solar Flare Forecasting Using Machine Learning and SDO/HMI Data: A Multiple Machine Learning Model and Data Curation Technique Comparison Study,'' \textit{The Astrophysical Journal Supplement Series}, vol.~280, no.~2, art.~54, 2025.
  \item \me, W.~D.~Pan, A.~Prasad, G.~A.~Gary, and Q.~Hu, ``On the Estimation of the SHARP Parameter MEANALP from AIA Images Using Deep Neural Networks,'' \textit{Solar Physics}, vol.~296, no.~11, art.~163, 2021.
\end{pubs}

\vspace{0.6ex}
\textbf{\itshape Conference and Workshop Papers, Abstracts, and Posters}
\begin{pubs}
  \item S.~Bonasera, G.~Acciarini, J.~A.~P\'erez-Hern\'andez, \me, E.~Brown, \textit{et al.}, ``Dropout and Ensemble Networks for Thermospheric Density Uncertainty Estimation,'' \textit{Bayesian Deep Learning Workshop, NeurIPS 2021}, 2021.
  \item E.~J.~E.~Brown, S.~Bonasera, \me, J.~A.~P\'erez-Hern\'andez, G.~Acciarini, \textit{et al.}, ``Learning the Solar Latent Space: Sigma-Variational Autoencoders for Multiple Channel Solar Imaging,'' \textit{Fourth Workshop on Machine Learning and the Physical Sciences, NeurIPS 2021}, 2021.
  \item \me, E.~Brown, S.~Bonasera, G.~Acciarini, J.~A.~P\'erez-Hern\'andez, \textit{et al.}, ``Simultaneous Multivariate Forecast of Space Weather Indices Using Deep Neural Network Ensembles,'' \textit{Fourth Workshop on Machine Learning and the Physical Sciences, NeurIPS 2021}, 2021.
  \item \me, Z.~Jiang, W.~D.~Pan, G.~A.~Gary, and Q.~Hu, ``Determination of Linear Force-Free Magnetic Field Constant Alpha Using Deep Learning,'' \textit{International Conference on Computational Science and Computational Intelligence (CSCI)}, 2017.
  \item T.~Singh, L.~Farris, C.~Hall, T.~Newman, \me, S.~Raza, and N.~Pogorelov, ``Solar Flare Forecasting Using Multiple Machine Learning Models and SDO/HMI Data,'' \textit{SHINE 2024 Workshop}, 2024.
  \item R.~Hooda, W.~D.~Pan, and \me, ``Transform Decomposition Switching for Efficient Attribute Compression of 3D Point Clouds Using Neural Networks,'' \textit{International Conference on Computational Science and Computational Intelligence (CSCI)}, 2022.
  \item T.~Singh, \me, \textit{et al.}, ``Improving the Arrival-Time Prediction of Coronal Mass Ejections Using Magnetohydrodynamic Ensemble Modeling, Heliospheric Imager Data, and Machine Learning,'' \textit{Proceedings of the 2nd Machine Learning in Heliophysics Conference}, 2022.
  \item C.~Davis, B.~Praveen, E.~Foshee, K.~Bugbee, D.~Sharma, N.~Pantha, S.~Awale, \me, \textit{et al.}, ``LLMs, Encoders, and Agentic Workflows for Curation and Discovery in NASA's Science Discovery Engine,'' \textit{AGU Fall Meeting 2025}, 2025.
  \item T.~Singh, C.~Hall, T.~Newman, \me, S.~Raza, and N.~Pogorelov, ``Solar Flare Forecasting Using Machine Learning and SDO/HMI Data: A Comparison of Multiple ML Models and AR Parameter Time Series,'' \textit{SHINE 2023 Workshop}, 2023.
  \item T.~Singh, \me, \textit{et al.}, ``Forecasting Solar Flares Using the Time Evolution of Active Regions and Machine Learning Techniques,'' \textit{NASA Proposal ID 21-SWR2O2R\_2-0011}, 2021.
\end{pubs}

% ----------------------------- INVITED TALKS --------------------------------
\section{Invited Talks}
\begin{itemize}
  \item ``Deep Neural Networks Applied to Three Problems in Heliophysics: Coronal Magnetic Field Modeling, Free Magnetic Energy, and Forecasting Solar Activity,'' NASA Marshall Space Flight Center, Heliophysics Division, Sep.~2020.
  \item ``Applications of AI in Logistics and Supply Chain,'' Auburn University, Harbert College of Business, Oct.~2023.
  \item ``AI in the Supply Chain,'' Mississippi State University, College of Business, Mar.~2024.
\end{itemize}

% ------------------------------- TEACHING -----------------------------------
\section{Teaching Experience}
\begin{itemize}
  \item \textbf{CPE211L --- Introduction to Programming (C++)}, undergraduate lab instructor, University of Alabama in Huntsville.
  \item \textbf{EE316 --- Electrical Circuits and Electrical Design Lab}, undergraduate lab instructor, University of Alabama in Huntsville (Summer 2017).
\end{itemize}

% --------------------------- HONORS & AWARDS --------------------------------
\section{Honors \& Awards}
\begin{itemize}
  \item Team Excellence Award, Frontier Development Lab (NASA / SETI Institute), 2021.
\end{itemize}

% --------------------------- PROFESSIONAL SERVICE ---------------------------
\section{Professional Service}
\textbf{Reviewer}
\begin{itemize}
  \item Research proposals, NASA Heliophysics.
  \item \textit{Journal of Geophysical Research: Machine Learning and Computation.}
  \item \textit{Journal of Atmospheric and Solar-Terrestrial Physics.}
  \item \textit{Solar Physics.}
  \item Journals of the American Astronomical Society.
  \item IEEE SoutheastCon --- AI and ML tracks (Mar.~2021).
\end{itemize}
\textbf{Session Chair}
\begin{itemize}
  \item IEEE SoutheastCon --- Computing track (Mar.~2021).
\end{itemize}

% ------------------------ SELECTED RESEARCH PROJECTS ------------------------
\section{Selected Research Projects}
\begin{itemize}
  \item \textbf{Improving CME Arrival Time with Ensemble Models and Machine Learning} --- regression and neural-network models to improve coronal mass ejection arrival-time predictions.
  \item \textbf{Simultaneous Multivariate Forecasts of Space-Weather Indices using DNN Ensembles} --- LSTM ensembles forecasting space-weather indices over a full solar rotation.
  \item \textbf{Solar Image Processing for Flare Detection} --- spatial and frequency-domain filtering to enhance coronal loops, latent-space visualization with VAEs, and relating SHARP parameters to flaring events.
  \item \textbf{Time-Series Forecast of Solar Cycle 25 using Deep Neural Networks} --- WaveNet and LSTM models predicting the strength of Solar Cycle 25.
\end{itemize}

% ----------------------------- TECHNICAL SKILLS -----------------------------
\section{Technical Skills}
\begin{itemize}
  \item \textbf{Programming Languages:} Python, SQL, C++, MATLAB.
  \item \textbf{Machine Learning \& Deep Learning:} PyTorch, TensorFlow, Keras, scikit-learn; CNNs, RNNs, LSTMs, Transformers; retrieval-augmented generation (RAG) and large language models.
  \item \textbf{Data Analysis \& Visualization:} Pandas, NumPy, SciPy, Matplotlib, Seaborn, Plotly, Power BI, Tableau.
  \item \textbf{Cloud \& DevOps:} AWS, Azure, Google Cloud Platform; Docker, Git, CI/CD pipelines, REST APIs.
\end{itemize}

\end{document}
